% 連続体力学 連続体力学-1 問題6（006-incomp） 解いた日 2026-10-21
%% kotae
\toi{1}{◎}{}
\[ \pdv{\rho}{t}=-\rho_{0}\varepsilon\kappa U\cos(\kappa(x_{1}-Ut)) \]
\[ \pdv{}{x_{1}}(\rho u_{1})=\rho_{0}\varepsilon\kappa U\cos(\kappa(x_{1}-Ut)) \]
よって
\[ \pdv{\rho}{t}+\vect{\nabla}\cdot(\rho\vect{u})=0 \]
\toi{2}{◎}{}
$\vect{\nabla}\cdot\vect{u}=0$　非圧縮
\[ \frac{\mathrm{D}\rho}{\mathrm{D}t}=\pdv{\rho}{t}+u_{1}\pdv{\rho}{x_{1}}=0 \]
（矢印で上の式から次の文につなぐ）

これが示しているのは、流れている粒子の密度が一定ということであって、密度が一様でないこととは矛盾しない。
\toi{3}{◎}{}
\[ \pdv{\rho}{t}=0 \]
\[ \pdv{}{x_{1}}(\rho u_{1})=\frac{\gamma\rho_{0}x_{2}}{L} \]
\[ \pdv{\rho}{t}+\vect{\nabla}\cdot(\rho\vect{u})\neq0 \]
より、保たれない。
\toi{4}{△}{計算}
$\vect{\nabla}\cdot\vect{u}=0$ より非圧縮
\[ \frac{\mathrm{D}\rho}{\mathrm{D}t}=\pdv{\rho}{t}+u_{1}\pdv{\rho}{x_{1}}+0=0 \]
\[ \therefore\ \pdv{\rho}{t}=-\frac{\gamma\rho_{0}x_{2}}{L} \]
\naze{ここで $\partial\rho/\partial x_{1}$ を時刻0の値 $\rho_{0}/L$ のまま使っている。時間がたっても変わらない理由が要る。粒子の軌道 $x_{2}=X_{2}$，$x_{1}=X_{1}+\gamma X_{2}t$ から出発点 $X_{1}=x_{1}-\gamma x_{2}t$ を出して初期の密度に代入するか，得た $\rho$ を代入して $\partial\rho/\partial x_{1}=\rho_{0}/L$ のままであることを確かめれば閉じる。}
\[ \rho(\vect{x},t)=-\frac{\gamma\rho_{0}}{L}x_{2}t+C \]
\[ \rho(\vect{x},0)=\rho_{0}(1+x_{1}/L)=C \]
\[ \rho(\vect{x},t)=\rho_{0}\Bigl(1+\frac{x_{1}}{L}-\frac{\gamma x_{2}}{L}t\Bigr) \]
\[ x_{1}-\gamma x_{2}t=K\quad (K\text{ は定数}) \]
（図：$x_{1}x_{2}$ 平面に，原点の近くを通る傾いた直線が2本。直線の上端の間に回転の向きを示す矢印）

\ayamari{反時計回り}{時計回り}に回転する。
\naze{$t=0$ では $x_{1}=K$（$x_{1}$ 軸に垂直）で，時刻 $t$ では $x_{1}=K+\gamma t\,x_{2}$ である。$x_{2}>0$ の部分は $+x_{1}$ 向きに流されるので，等密度面は $x_{2}$ 軸から角度 $\phi$（$\tan\phi=\gamma t$）だけ時計回りに倒れ，流れの向き（$x_{1}$ 軸）に近づいていく。図の直線の傾きは正しいが，矢印と文が逆向きである。}
%% kizuki
\item 等密度面の回転の向きは，1点を流れに沿って動かして確かめる。$x_{2}>0$ の点は $+x_{1}$ 向き，$x_{2}<0$ の点は $-x_{1}$ 向きに流されるので，$x_{2}$ 軸に沿った線は上が右へ，下が左へずれて時計回りに倒れる（$x_{1}$ 軸に近づく）。
\item \mae{第1回 問題6・10-17}前回は $\mathrm{D}\rho/\mathrm{D}t=0$ で止まった。今回は密度場まで出せたが，$\partial\rho/\partial x_{1}$ を一定として $\partial\rho/\partial t$ を出す所に根拠が無い。$\mathrm{D}\rho/\mathrm{D}t=0$ は「各粒子が時刻0の密度を運ぶ」ことなので，軌道から出発点 $X_{1}=x_{1}-\gamma x_{2}t$ を求めて初期の密度に代入するのが確実で，得た $\rho$ は $\partial\rho/\partial t+\gamma x_{2}\,\partial\rho/\partial x_{1}=0$ に代入して検算する。
